\begin{ex}
Để vẽ đoạn thẳng AB dài 7 cm, ta đặt thước kẻ sao cho điểm A trùng với vạch số mấy của thước?
\shortans{0}
\loigiai{
Điểm A trùng vạch 0 cm.
}
\end{ex}

\begin{ex}
Điểm B sẽ trùng với vạch số mấy của thước kẻ?
\shortans{7}
\loigiai{
Điểm B trùng vạch số 7 cm.
}
\end{ex}

\begin{ex}
Cắt một tờ giấy hình vuông thành hai hình tam giác bằng nhau bởi một đường chéo. Cần bao nhiêu nhát kéo cắt thẳng?
\shortans{1}
\loigiai{
Cần 1 nhát kéo cắt theo đường chéo.
}
\end{ex}

\begin{ex}
Gấp đôi một tờ giấy hình chữ nhật theo chiều ngang, ta nhận được 2 hình gì bằng nhau? Mỗi phần là hình chữ nhật nhỏ. Có bao nhiêu hình chữ nhật nhỏ?
\shortans{2}
\loigiai{
Có 2 hình chữ nhật nhỏ.
}
\end{ex}

\begin{ex}
Dùng 4 que tính bằng nhau xếp được 1 hình vuông. Muốn xếp 3 hình vuông rời nhau cần bao nhiêu que tính?
\shortans{12}
\loigiai{
Cần: $4 \times 3 = 12$ que tính.
}
\end{ex}

\begin{ex}
Dùng 3 que tính bằng nhau xếp được 1 hình tam giác. Muốn xếp 4 hình tam giác rời nhau cần bao nhiêu que tính?
\shortans{12}
\loigiai{
Cần: $3 \times 4 = 12$ que tính.
}
\end{ex}

\begin{ex}
Dùng 7 que tính bằng nhau xếp được 2 hình vuông chung 1 cạnh. Cần bao nhiêu que tính?
\shortans{7}
\loigiai{
Cần đúng 7 que tính.
}
\end{ex}

\begin{ex}
Dùng 5 que tính bằng nhau xếp được 2 hình tam giác chung 1 cạnh. Cần bao nhiêu que tính?
\shortans{5}
\loigiai{
Cần đúng 5 que tính.
}
\end{ex}

\begin{ex}
Một miếng bìa hình vuông có chu vi 16 cm. Cạnh của miếng bìa hình vuông dài bao nhiêu xăng-ti-mét?
\shortans{4}
\loigiai{
Cạnh dài: $16 : 4 = 4$ cm.
}
\end{ex}

\begin{ex}
Ghép 2 hình tam giác vuông cân bằng nhau có thể tạo thành hình gì có 4 cạnh bằng nhau? Hình vuông đó có bao nhiêu cạnh?
\shortans{4}
\loigiai{
Tạo thành hình vuông có 4 cạnh.
}
\end{ex}

\begin{ex}
Ghép 6 hình tam giác đều nhỏ bằng nhau có thể tạo thành một hình lục giác đều. Hình lục giác đều có bao nhiêu cạnh?
\shortans{6}
\loigiai{
Hình lục giác có 6 cạnh.
}
\end{ex}

\begin{ex}
Cắt bỏ 1 góc của hình vuông bằng 1 đường cắt thẳng. Hình mới nhận được có bao nhiêu cạnh?
\shortans{5}
\loigiai{
Hình mới có 5 cạnh.
}
\end{ex}

\begin{ex}
Cắt đôi một hình vuông bằng một đường thẳng đi qua trung điểm hai cạnh đối diện, ta nhận được bao nhiêu hình chữ nhật?
\shortans{2}
\loigiai{
Nhận được 2 hình chữ nhật.
}
\end{ex}

\begin{ex}
Một hình chữ nhật được ghép từ 3 hình vuông nhỏ cạnh 2 cm xếp thành hàng ngang. Chiều dài hình chữ nhật đó là bao nhiêu xăng-ti-mét?
\shortans{6}
\loigiai{
Chiều dài: $2 \times 3 = 6$ cm.
}
\end{ex}

\begin{ex}
Chiều rộng của hình chữ nhật đó là bao nhiêu xăng-ti-mét?
\shortans{2}
\loigiai{
Chiều rộng bằng cạnh hình vuông là 2 cm.
}
\end{ex}

\begin{ex}
Chu vi của hình chữ nhật ghép từ 3 hình vuông đó là bao nhiêu xăng-ti-mét?
\shortans{16}
\loigiai{
Chu vi: $(6 + 2) \times 2 = 16$ cm.
}
\end{ex}

\begin{ex}
Có 8 que tính dài bằng nhau, xếp được nhiều nhất bao nhiêu hình vuông rời nhau?
\shortans{2}
\loigiai{
Xếp được: $8 : 4 = 2$ hình vuông.
}
\end{ex}

\begin{ex}
Có 9 que tính dài bằng nhau, xếp được nhiều nhất bao nhiêu hình tam giác rời nhau?
\shortans{3}
\loigiai{
Xếp được: $9 : 3 = 3$ hình tam giác.
}
\end{ex}

\begin{ex}
Một tờ giấy hình chữ nhật có chiều dài 15 cm, chiều rộng 10 cm. Cắt một hình vuông có cạnh 10 cm từ tờ giấy đó. Mảnh giấy còn lại là hình chữ nhật có chiều dài 10 cm và chiều rộng bao nhiêu xăng-ti-mét?
\shortans{5}
\loigiai{
Chiều rộng mảnh còn lại: $15 - 10 = 5$ cm.
}
\end{ex}

\begin{ex}
Bác thợ mộc có 10 thanh gỗ dài bằng nhau để đóng các khung tranh hình tam giác rời nhau. Bác đóng được nhiều nhất bao nhiêu khung tranh?
\shortans{3}
\loigiai{
Vì $10 : 3 = 3$ (dư 1) nên đóng được nhiều nhất 3 khung tranh.
}
\end{ex}

\begin{ex}
Lan dùng các que tính xếp thành chữ số 0 trên bảng điện tử. Chữ số 0 cần bao nhiêu que tính?
\shortans{6}
\loigiai{
Cần 6 đoạn que tính (2 ngang, 4 dọc).
}
\end{ex}

\begin{ex}
Chữ số 8 trên bảng điện tử cần bao nhiêu que tính?
\shortans{7}
\loigiai{
Cần 7 đoạn que tính.
}
\end{ex}

\begin{ex}
Chữ số 7 trên bảng điện tử cần bao nhiêu que tính?
\shortans{3}
\loigiai{
Cần 3 đoạn que tính.
}
\end{ex}

\begin{ex}
Xếp một hình ngôi nhà gồm 1 hình tam giác làm mái và 1 hình vuông làm thân. Ngôi nhà có tất cả bao nhiêu cạnh bên ngoài bao quanh?
\shortans{5}
\loigiai{
Bao quanh gồm 2 cạnh mái và 3 cạnh thân nhà, tổng là 5 cạnh.
}
\end{ex}

\begin{ex}
Để vẽ đoạn thẳng CD dài 14 cm, điểm C ở vạch 2 cm thì điểm D ở vạch số mấy của thước kẻ?
\shortans{16}
\loigiai{
Điểm D ở vạch: $2 + 14 = 16$ cm.
}
\end{ex}
